% 生成AI等を利用していない場合は，このファイルをコメントだけのままにします．
% コメントだけなら見出しも空白頁も出力されません．
%
% 生成AI等を利用した場合は，提出時点の大学・学部・学科の規則と
% 指導教員の指示を先に確認し，実際の利用事実だけを次の形式で記載します．
% 必要事項は，サービス名・モデル名・版または利用期間，用途と利用範囲，
% 入力した情報の範囲，出力を検証した対象と方法，著者の責任です．
% 入出力や作業履歴の保存を求められた場合は，指定された方法で保管します．
% 生成AI等がデータ処理，解析，判定，評価指標の算出等の研究方法に
% 関与した場合は，この申告だけで済ませず，方法・実験の章にも再現可能な
% 条件，検証方法および限界を記載します．
% 内容入りの架空例はexamples/ai-usage-example.texにあります．
%
% \chapter*{生成AI等の利用に関する申告}
% \markboth{生成AI等の利用に関する申告}{生成AI等の利用に関する申告}
% \addcontentsline{toc}{chapter}{生成AI等の利用に関する申告}
%
% 本論文の作成にあたり，著者は［サービス名・モデル名・版または利用期間］を，
% ［実際の用途］に使用した．利用範囲は［対象とした章，文章，コード等］である．
% ［入力した情報の範囲と，外部サービスへ入力しなかった情報を記載する］．
% 生成AI等の出力について，著者は［原資料，一次文献，実行結果等］と照合し，
% ［検証した対象と方法］を確認した．生成AI等を著者または根拠資料として扱わず，
% 本論文の正確性，研究公正，独創性および最終内容については著者が責任を負う．
